Magnetic-core models used in power-converter design are commonly fitted
to nominal material curves despite limited operating-point coverage and
material variation. We present a Bayesian framework that calibrates
Steinmetz core-loss and Cole--Cole complex-permeability models and
evaluates expected information gain (EIG) for sequential measurement
selection. Five prior-predictive, matched-model recovery cases gave
per-parameter median absolute relative errors from 0.27\% to 5.85\%,
with the generating values inside 28 of 30 equal-tailed 90\%
model-conditional credible intervals. In a separate 30-paired-seed
benchmark with shared candidate outcomes, raw EIG reached a prespecified
two-target latent-response precision gate in five measurements, compared
with nine for a deterministic fixed channel-balanced traversal. Against
stronger comparators, however, raw EIG tied predictive variance and
Laplace D-optimality on measurement count in all 30 pairs. EIG per
modeled cost tied Laplace D-optimality but required a mean 15.17 more
modeled-cost units than predictive variance (95\% paired-bootstrap
interval 15.00--15.50 units in favor of predictive variance). No policy
failed the gate. On accepted public measured records, in-sample RRMSE
was 8.79\%--18.21\% for core loss, 6.89\%--9.33\% for \(\mu'\), and
36.77\%--52.42\% for \(\mu''\); the loss-component residuals expose
substantial one-pole model discrepancy. The evidence supports a
reproducible, model-conditional acquisition benchmark. A separate
preregistered 120-task synthetic mismatch campaign found that raw EIG
reached the same local gate in every combined-mismatch seed but was
false-confident in 22 of 30; its strong-comparator count advantages were
only 0.10 measurements. This shows that local posterior width did not
protect truth accuracy under the locked departures. It does not show EIG
superiority over strong comparators, measured laboratory-time savings,
global six-parameter identification, or a validated optimal laboratory
plan.
\href{https://github.com/hoangduong6210/EIG-bayesian-for-Recover-potential-Physical-Parameter-of-MagComponent/blob/927bbcb6d0b862598fae87d2575a872bbd79b5f1/wiki/evidence/Evidence-Sources.md\#e2}{Evidence
E2},
\href{https://github.com/hoangduong6210/EIG-bayesian-for-Recover-potential-Physical-Parameter-of-MagComponent/blob/927bbcb6d0b862598fae87d2575a872bbd79b5f1/wiki/evidence/Evidence-Sources.md\#e4}{E4},
\href{https://github.com/hoangduong6210/EIG-bayesian-for-Recover-potential-Physical-Parameter-of-MagComponent/blob/927bbcb6d0b862598fae87d2575a872bbd79b5f1/wiki/evidence/Evidence-Sources.md\#e7}{E7},
and
\href{https://github.com/hoangduong6210/EIG-bayesian-for-Recover-potential-Physical-Parameter-of-MagComponent/blob/927bbcb6d0b862598fae87d2575a872bbd79b5f1/wiki/evidence/Evidence-Sources.md\#e16}{E16}
